\documentclass[11pt]{article}
\usepackage[T1]{fontenc}
\usepackage[margin=1in]{geometry}
\usepackage{microtype}
\usepackage[hidelinks]{hyperref}
\hypersetup{pdftitle={Thoughts and Notes on Energy-based Neuron and Intelligence}}

\setlength{\parindent}{0pt}
\setlength{\parskip}{0.65em}
\pagestyle{plain}

\begin{document}

\begin{center}
{\Large Thoughts and Notes on Energy-based Neuron and Intelligence}
\end{center}

\vspace{0.5em}

\section*{In search of a better neuron}

What if a novel artificial neuron with more leverage in the system is all you need? The dominant paradigm in deep learning use on backpropagation flow through layers to update weights (globally) as it has withstand the test of time and scaling. Yet biological neurons learn through fundamentally different mechanisms (at least we think they do), relying on local signals and activity-dependent plasticity probably without any global error gradient. This difference between artificial and biological neural nets suggests an alternative path worth exploring. Can we build neural networks that learn like brains do, using only local information and achieve meaningful results? I pursued this idea through several experiments, developing architectures where neurons learned based on local states.

The core (gut) idea is a recursive neuron. A recursive neuron adds additonal gate after activation that accumulates energy and decides when to ``fire'' based on that internal state. They have internal memory that persists across time, and they decide when to update that memory and when to hold onto previous state, decision not done by external mechanism or learned through backpropagation; it emerges from the neuron's dynamics.

Neuron gets input $x$, projects it through weights $W$ to compute drive $d = Wx$. Before we add to memory: $\tilde{d} = d / \sigma$ where $\sigma$ is running variance. This stops large inputs from overwhelming memory. Memory updates as $m_t = \alpha m_{t-1} + \tilde{d}$, an exponential moving average (EMA). Then neuron computes a residual ratio $r = \|m_t - m_{t-1}\| / \|\tilde{d}\|$ measuring how much the new input changed memory. If $r >$ threshold, the neuron fires and outputs $m_t$. Otherwise it propagates 0 and keeps accumulating.

I tried this approach on synthetic temporal tasks---patterns where the model learns structure, and delay-query tasks where it must hold information across time before giving output. Results: entropy $< 0.1$, hit rates $> 0.95$ across five random seeds. Learning utilized energy-gated Hebbian plasticity as $\Delta W \propto E \cdot xy$ where $E$ is accumulated energy, $x$ is pre-synaptic activity, $y$ is post-synaptic activity.

Scaling to real data was harder. I built a language model with multiple recursive neuron layers. First layer processes character embeddings, later layers process working memory from below.

On War and Peace, the model hit cross-entropy 3.09 after 80k steps. Random guessing gives 4.4, transformers get \textasciitilde{}1.5, theoretical floor is \textasciitilde{}1.3. The model learned some structure but couldn't learn representations end-to-end. Each layer optimizes locally, not for the global objective - no gradient coordination means a performance ceiling.

\section*{Intelligence}

The definition of intelligence is more a game of words than numbers. I do believe intelligence is quantifiable. Do we have the correct tools to do so? No. Not for humans and neither for machines (AI systems).  I do not believe general intelligence is quantifiable. I like suprisingly satifsying non-rigorous definition of intelligence:

\begin{quote}
Intelligence is a process of local minimization of entropy in the known universe.
\end{quote}

Or rephrased poetically (inspired by Yevgeny Zamyatina famous qoute in We):

\begin{quote}
Intelligence is the only bitter (local) remedy against the entropy of the universe.
\end{quote}

which is a technically incorrect, as there are other ways of locally minizing entropy - or at least intepretations - yet I wrote down this one as it is sound. For context I will list some of the popular quotes on intelligence stated by achieved scholars:

\textbf{Alan Turing} (computing pioneer):
\begin{quote}
``The original question, `Can machines think?' I believe to be too meaningless to deserve discussion.''
\end{quote}

\textbf{Marvin Minsky} (AI researcher):
\begin{quote}
``Intelligence is a suitcase word. You can't define it until you open the suitcase and see what's inside.''
\end{quote}

\textbf{Jeff Hawkins} (neuroscientist, Numenta):
\begin{quote}
``Intelligence is measured by the capacity to predict.''
\end{quote}

\textbf{Karl Friston} (neuroscientist, free energy principle):
\begin{quote}
``Intelligent behavior is simply inference about the future that minimizes surprise.''
\end{quote}

\textbf{David Deutsch} (physicist, quantum computing):
\begin{quote}
``Intelligence is the ability to create new explanations.''
\end{quote}

\textbf{Douglas Hofstadter} (cognitive scientist):
\begin{quote}
``Intelligence is the ability to take a broad view, to get the gist of things, to see analogies.''
\end{quote}

These definitions cluster around prediction, adaptation, goal-achievement, and explanation while being underlied by compression, efficiency, and the ability to generalize from the particular to the general. By these terms majority of human population is at least subconciously intelligent.

\end{document}
